\documentclass[aspectratio=169]{beamer}
\usetheme{metropolis}
\metroset{sectionpage=none, numbering=counter, progressbar=frametitle}
\usepackage{amsmath, amssymb}
\usepackage[T1]{fontenc}

\title{Domain-Partition Algorithm for Turbomachinery Blades}
\subtitle{Slide 6: Singularities --- Index \& Poincar\'{e}-Hopf}
\date{\today}
\author{}
\institute{}

\begin{document}

\begin{frame}{Singularities --- Index \& Poincar\'{e}-Hopf}

\begin{center}
\Large
\[
\text{index}(p) = \frac{1}{2\pi} \oint_{\partial B(p)} d\theta
\]
\end{center}

\vspace{0.2em}
\begin{itemize}
    \item \textbf{4 tip singularities,} each with index $-\tfrac{1}{4}$
    \item \textbf{Poincar\'{e}-Hopf:} sum of indices equals Euler characteristic $\chi$
    \item \textbf{Periodic cylinder:} $\chi = 0$, so indices must balance
\end{itemize}

\vspace{0.2em}
\begin{center}
\[
\sum_i \text{index}(p_i) = \chi(\Omega)
\]
\end{center}

\vspace{0.1em}
\small
\textit{Detection:} \texttt{tools/singularity\_detector.py} \quad
\textit{Lit:} Kowalski 2015, ``Blocking piecewise-linear domains''

\note{
Speaker notes:
Singularities are points where the frame field is undefined. The index measures how much the field rotates around the point, computed as (1/2pi) times the integral of d theta around a small loop. For the T1_9 turbine hub we find exactly 4 tip singularities, each carrying index negative one quarter. The Poincare-Hopf theorem demands that the sum of all indices equals the Euler characteristic of the surface. Because the unwrapped hub is a periodic cylinder, chi is 0, so the four negative quarter indices sum to negative 1 and must be compensated by other singularities or boundary contributions. The singularity detector in tools/singularity_detector.py performs this classification automatically.
}

\end{frame}

\end{document}
